% !TEX program = xelatex
\documentclass{ctexart}
\usepackage{xeCJK}
\usepackage{fontspec}
\usepackage{graphicx} % Required for inserting images
\usepackage{amsmath} % Required for some math elements
\usepackage{pdfpages}
\usepackage{fancyhdr}
\usepackage{tocloft}
\usepackage{setspace}
\usepackage[backend=biber,style=gb7714-2015, url=false]{biblatex}  
\usepackage{etoolbox} % 提供 \patchcmd 命令
\usepackage[colorlinks=true, linkcolor=black, citecolor=black]{hyperref}

\newcommand{\simsun}{\CJKfamily{SimSun}}
\newcommand{\simhei}{\CJKfamily{SimHei}}
\newcommand{\simkai}{\CJKfamily{KaiTi}}
\newcommand{\STSimsun}{\CJKfamily{STZhongsong}}

\setCJKmainfont{SimSun}
\setmainfont{Times New Roman}
\addbibresource{ref.bib}

% ===== 一级标题左对齐 =====
\ctexset{
  section/format = \raggedright\bfseries\zihao{3},
}



\pagestyle{empty}
\fancyhf{}

% ===== 目录专用设置 =====
% 1. 行距设置 (仅作用于目录)
\let\oldtoc\tableofcontents % 保存原目录命令
\renewcommand{\tableofcontents}{% 重新定义目录命令
  \begin{spacing}{1.5} % 1.5倍行距
    \oldtoc% 调用原目录命令
  \end{spacing}
}

% 2. 字体设置 (小四号宋体)
\renewcommand{\cfttoctitlefont}{\normalfont\CJKfamily{SimSun}\huge} % 目录标题字体
\renewcommand{\cftsecfont}{\normalfont\CJKfamily{SimSun}} % 节条目字体
\renewcommand{\cftsubsecfont}{\normalfont\CJKfamily{SimSun}} % 子节条目字体

% 3. 调整缩进
\setlength{\cftsecindent}{0em} % 节标题缩进
\setlength{\cftsubsecindent}{2em} % 子节缩进

\begin{document}

\setmainfont{Times New Roman}
% \setCJKmainfont{SimSun}

% 封面 Covers/cover.pdf

% ===== 重定义标题格式 =====
\newfontfamily\titlelatinfont[Scale=MatchUppercase]{Times New Roman}
\makeatletter
\renewcommand\@maketitle{
    \begin{center}
        {\zihao{2}\STSimsun\titlelatinfont
        \@title\par}
        \vskip 1em
        {\large\@author}
    \end{center}
}
\makeatother


\newenvironment{zhabstract}{%
  \par\noindent
  {\zihao{-4}\normalfont\simkai\textbf{【摘要】}} % 标题：小四号楷体加粗
  \par\vspace{0.5em}
}
{\par\vspace{0.5em}}


% ===== 关键词环境定义 =====
\newenvironment{zhkeywords}{%
  \par\noindent
  {\zihao{-4}\normalfont\simkai\textbf{【关键词】}} % 标题：小四号楷体加粗
  \par\vspace{0.5em}
}
{\par\par\vspace{0.5em}}


\newenvironment{acknowledgements}{%
  \vspace{0.5\baselineskip} % 段前 0.5 行
  \par\noindent
  {\zihao{-4}\normalfont\heiti\textbf{致谢}} % 标题：小四号黑体加粗
  \par\vspace{0.5\baselineskip} % 段后 0.5 行
}

\newpage
\tableofcontents


\newpage
\setcounter{page}{1}
% Title
\title{一种基于STAGE分离和异构部署的多模态推理优化系统}
\author{高忠山、曾立}
\maketitle
\begin{zhabstract}
-
\end{zhabstract}
\begin{zhkeywords}
-
\end{zhkeywords}

% ===== 1. 引言 =====
\section{引言}

多模态大模型在图像生成领域的主要任务有文生图（T2I），图生图（I2I）,图文编辑（TI2I），其规模化应用对推理系统提出了高并发、低延迟的严峻挑战。当前主流的多模态图像生成模型如Qwen-Image~\cite{QwenImageTechnicalReportWu2025}采用VAE编码器（I2I，TI2I）-大语言模型特征编码-MMDiT扩散去噪-VAE解码器的复合架构，推理流程涉及参数规模、计算特征迥异的多个环节。然而，现有推理服务系统如vLLM-Omni~\cite{VLLMOmniFullyDisaggregatedYin2026}虽已实现对自回归多模态模型的stage分离推理，
但普遍未实现stage深度分离与异构硬件精准匹配，存在显著技术缺陷，无法满足高并发、低延迟的规模化推理需求。由此导致三大核心问题：（1）异构硬件资源利用率低，不同算力等级的AI处理器与CPU间协同效率低下；（2）推理各阶段耦合度高，资源争抢严重，端到端延迟难以控制；（3）大模型权重加载耗时过长，冷启动延迟无法满足高频次推理服务需求。

针对上述问题，本文提出一种基于stage分离和异构部署的多模态推理优化系统。主要贡献如下：
\begin{enumerate}
    \item 将多模态模型的文生图推理pipeline拆分为文本预处理、LLM特征编码、MMDiT Diffusion生成、VAE Decoder四个功能独立、边界清晰的stage，并针对昇腾NPU多型号（310P/910C/910B 等）+ CPU的异构硬件环境，建立基于算力特征、参数规模、耗时占比的硬件亲和性匹配规则，实现各stage与异构硬件的精准适配。
    \item 设计跨硬件异步流水线调度机制，为每个stage维护专属任务队列并结合动态批处理策略，消除阶段间同步等待，显著提升异构硬件并行效率与推理吞吐量。
    \item 构建适配stage分离架构的三级分层缓存（模型权重缓存、中间特征缓存、推理参数缓存），实现阶段级缓存预加载与淘汰更新，大幅缩短模型冷启动时间。
\end{enumerate}

% ===== 2. 相关工作 =====
\section{相关工作}

\subsection{大模型与多模态推理服务}

以大语言模型（LLM）高效推理服务为核心，vLLM~\cite{EfficientMemoryManagementKwon2023}提出PagedAttention机制实现KV-cache内存近零浪费管理，显著提升LLM推理吞吐量。Miao等~\cite{EfficientGenerativeLargeMiao2026}从算法创新与系统优化两个维度系统综述了生成式LLM推理服务的效率优化方法，覆盖解码算法、模型压缩、并行计算、内存管理、请求调度等关键方向。
SARATHI~\cite{SARATHIEfficientLLMAgrawal2023}通过chunked-prefill与decode-maximal batching策略，将计算密集的prefill与内存密集的decode阶段统一调度，有效消除流水线气泡。
在多模态服务方向，
vLLM-Omni~\cite{VLLMOmniFullyDisaggregatedYin2026}提出面向Any-to-Any多模态模型的完全分离式服务架构。
对Thinker-Talker（如Qwen2.5-Omni,Qwen3-Omni等）、AR+DiT（如GLM-Image等），AR+专用生成器（如BAGEL等），通过stage级别抽象实现跨模态组件的独立调度与资源分配。

\subsection{分离式推理}

分离式推理（Disaggregated Inference）将LLM推理的prefill（计算密集）与decode（内存密集）两阶段拆分至不同硬件，实现独立资源分配与并行度优化。Splitwise~\cite{SplitwiseEfficientGenerativePatel2024}率先论证了prefill/decode阶段分离的经济性与可行性，通过将两个阶段部署至不同GPU集群实现1.4倍吞吐提升与20\%成本降低。DistServe~\cite{DistServeDisaggregatingPrefillZhong2024}进一步将分离策略与SLO约束协同优化，消除prefill-decode干扰。
TetriInfer~\cite{InferenceInterferenceDisaggregateHu2024}提出两级调度+资源预测的分离式架构，降低混合负载下的推理干扰。
现有分离式推理工作对DiT文生图流水线（文本编码→扩散去噪→VAE解码）多阶段分离的异构推理部署涉及较少。本工作将异构分离式推理范式拓展至扩散模型领域，提出面向文生图流水线的四阶段深度分离方案。

\subsection{扩散模型并行推理}

扩散模型的高分辨率图像生成面临巨大计算开销。
DistriFusion~\cite{DistriFusionDistributedParallelLi2024}提出displaced patch parallelism，利用相邻扩散步骤特征图的高相似性，通过复用前序步骤的stale feature map实现通信与计算异步重叠实现了最高6.1倍加速。
PipeFusion~\cite{PipeFusionPatchlevelPipelineFang2025}设计patch级流水线并行策略，将DiT模型层与图像patch分配至多GPU，复用一步stale特征降低通信开销，配合参数分布存储提升内存效率。xDiT~\cite{XDiTInferenceEngineFang2024}整合序列并行、PipeFusion流水线并行与CFG并行，构建可灵活组合的混合并行推理引擎，首次在以太网互联GPU集群上展示DiT可扩展性。DiT-Serve~\cite{DiTServeEfficientServingLuo2025}则从跨请求批处理优化角度出发，提出step-level batching，在每个denoising step边界抢占和交换请求以消除异构请求间的等待气泡，并设计Brick Attention将不同上下文长度的请求binpack到多组独立ring上，减少padding开销。
TridentServe~\cite{TridentServeStagelevelServingXia2025}针对扩散pipeline内在的stage异质性与请求异质性，提出动态stage级服务范式，通过自动placement plan和dispatch plan实现各stage的细粒度资源分配，消除静态部署导致的资源浪费。
上述并行推理工作主要聚焦于同构GPU环境下的模型内部加速或批处理调度优化，TridentServe虽实现了stage级动态分配但基于同构GPU集群，尚未覆盖异构加速器（如昇腾NPU多型号）的算力梯度差异。本工作将stage分离与异构硬件亲和性匹配相结合，面向昇腾310P/910C/910B等多型号NPU+CPU环境，提出四阶段深度分离与精准异构部署方案。
\subsection{异构部署}

异构计算环境下CPU与GPU/加速器的协同部署是降低推理成本的重要路径。HeteGen~\cite{HeteGenHeterogeneousParallelZhao2024}提出CPU+GPU异构张量并行框架，通过异步重叠通信与计算缓解I/O瓶颈，实现资源受限设备上的高效LLM推理。
SiPipe~\cite{SiPipeBridgingCPUGPUHe2025}利用未充分利用的CPU资源卸载辅助计算与通信任务，通过CPU sampling、token-safe执行模型与结构感知传输三项技术缓解流水线气泡，在多GPU流水线并行部署中提升GPU利用率达23\%。Hybrid SD~\cite{HybridSDEdgeCloudYan2024}将Stable Diffusion的早期扩散步骤部署于云端大模型实现语义规划，后期步骤部署于边缘端小模型完成视觉细节精炼，实现边缘-云协同扩散推理。现有异构部署方案止于CPU-GPU/云端-边缘的二元划分，尚未针对昇腾NPU多型号（310P/910C/910B 等）的算力梯度差异做逐stage的硬件亲和性精准匹配。本工作首次提出面向四阶段异构硬件的stage-硬件亲和性匹配方法。

\subsection{推理缓存技术}

缓存技术是缩短模型启动延迟、提升推理效率的关键手段。LMCache~\cite{LMCacheEfficientKVLiu2025}提出层次化KV缓存方案，支持GPU显存↔CPU内存↔磁盘↔远程存储的多层级缓存迁移与跨引擎复用，在prefix caching场景实现最高15倍吞吐提升。IMPRESS~\cite{IMPRESSImportanceInformedMultiTierChen2025}构建基于重要性的多级前缀KV存储系统，根据访问频率动态调整缓存层级，实现高命中率下的存储成本优化。PRESERVE通过模型权重与KV Cache的预取实现通信与计算重叠，缩短冷启动延迟。FlashPS~\cite{FlashPSEfficientGenerativeJiang2026}面向图像编辑场景，通过缓存和复用未编辑区域的中间activations跳过冗余计算，并设计bubble-free pipeline将缓存加载与计算流水线重叠，是扩散模型场景下结合缓存与计算调度的代表性工作。现有缓存方案几乎全部聚焦于KV-cache的层次化管理与跨请求复用，FlashPS虽涉及扩散模型的中间特征缓存但面向图像编辑的特殊mask场景，尚未出现适配通用文生图stage分离架构的模型权重+中间特征+推理参数三级缓存设计。本工作提出了与四阶段分离推理深度耦合的三级分层缓存机制，填补了这一空白。

% ===== 3. 系统设计 =====
\section{系统设计}

\subsection{系统整体架构}

本系统由四个核心模块组成：stage异构部署模块、跨硬件流水线调度模块、分层启动缓存模块、结果聚合模块。各模块数据互通、协同工作，形成闭环式推理优化架构。
\subsection{Stage分离与异构部署}

\subsubsection{Stage分离设计}
以Qwen-Image多模态大模型为例，针对多模态文生图推理任务场景，我们将模型推理流程拆分为4个边界清晰功能独立的stage，具体如下：
\begin{itemize}
  \item \textbf{文本预处理stage}：负责多模态中文本预处理与嵌入编码，将用户输入的文本提示词转换为初始特征向量，为后续文本特征深度提取提供基础；
  \item \textbf{LLM特征编码stage}：-负责对文本预处理stage输出的初始特征向量进行深度文本特征提取与编码，进一步优化文本信息向特征向量的转换效果；
  \item \textbf{MMDiT Diffusion生成stage}：-负责潜空间图像特征去噪过程，为多模态推理流程的计算密集型环节；
  \item \textbf{VAE解码stage}：-负责多模态中的图像重建，输出最终多模态文生图效果。
\end{itemize}
\subsubsection{异构硬件亲和性匹配}
基于多模态各stage的算力需求、参数规模、耗时占比、显存阈值，结合多模态双模态特性，精准分配至对应硬件，具体部署逻辑如下：
\begin{itemize}
  \item \textbf{文本预处理stage$\rightarrow$通用处理器（例如CPU）}：多模态中的轻量文本预处理与初始编码环节，例如参数规模约84M，耗时占比＜2\%，充分利用通用处理器资源，避免浪费AI处理器（例如昇腾NPU）算力，适配多模态轻量文本处理需求；
  \item \textbf{LLM特征编码stage$\rightarrow$中端AI处理器（例如昇腾NPU-310P）}：多模态中的轻量文本编码环节，例如参数规模7B，耗时占比＜3\%，适配中端AI处理器算力、24GB显存，释放高端AI处理器资源，适配多模态轻量文本处理需求；
  \item \textbf{MMDiT Diffusion生成stage$\rightarrow$高端AI处理器（例如昇腾NPU-910C）}：多模态图文双模态融合及核心生成环节，例如参数规模20B，耗时占比60\%-80\%，充分发挥高端AI处理器并行计算优势、64GB大显存特点，满足多模态密集型计算需求；
  \item \textbf{VAE解码stage$\rightarrow$对应AI处理器（例如昇腾NPU-910B）}：多模态中的图像重建环节，例如参数规模约100M，耗时占比10\%-20\%，适配32GB显存，与核心生成stage并行执行，提升多模态推理整体效率；
\end{itemize}
\subsubsection{设备拓扑}
-
\subsection{跨硬件流水线调度}

\subsubsection{任务队列管理}

系统为每个stage维护独立的输入队列与输出缓冲，形成四级解耦的任务队列架构：CPU(文本预处理)$\rightarrow$310P(LLM编码)$\rightarrow$910C(MMDiT去噪)$\rightarrow$910B(VAE解码)。队列条目包含请求ID、数据引用指针、中间特征引用及优先级标签。

为防止上下游处理速率失配导致内存膨胀，引入自适应背压机制：设各stage输入队列安全阈值$Q_{\text{safe}}$与上限$Q_{\text{max}}$，当队列深度$d$满足$Q_{\text{safe}}<d\le Q_{\text{max}}$时，上游分发间隔按比例$\frac{d-Q_{\text{safe}}}{Q_{\text{max}}-Q_{\text{safe}}}$递增；当$d>Q_{\text{max}}$时暂停投递。优先级调度采用两级策略：下游stage的输出队列优先级高于上游以加速端到端推进；短提示词优先于长提示词以减少排队延迟。

\subsubsection{动态批处理策略}

系统在MMDiT Diffusion stage上实现step级连续批处理：每个去噪步（denoising step）结束后，已完成全部步数的请求立即退出batch，其潜空间张量被转发至VAE解码队列，释放的显存被回收；同时输入队列中的新请求在一步内加入下一batch。该机制将batch空闲率从静态方案的约40\%降至5\%以下，与Orca~\cite{OrcaDistributedServingYu2022}在LLM推理中提出的连续批处理、DiT-Serve~\cite{DiTServeEfficientServingLuo2025}在DiT推理中提出的step-level batching思路一致。

各stage的批处理大小$B_s$根据三个约束实时计算并取最小值：$B_s^{\text{queue}}=\min(B_s^{\text{max}}, n_{\text{pending}})$（队列积压）、$B_s^{\text{mem}}=\lfloor (M_{\text{total}}-M_{\text{model}})/M_{\text{per\_req}} \rfloor$（显存约束）、$B_s^{\text{slo}}=\arg\max_k\{\text{Latency}(k)\le T_{\text{slo}}\}$（SLO约束）。各stage因硬件差异采用差异化上限：CPU文本预处理$B_{\text{text}}\le64$，310P端LLM编码$B_{\text{llm}}\le8$（24GB显存），910C端MMDiT$B_{\text{dit}}\le4$（1024$\times$1024单请求$\sim$16GB），910B端VAE$B_{\text{vae}}\le16$。

\subsubsection{异步流水线调度}

系统采用全异步流水线执行模型：stage完成当前任务后立即将输出写入下游输入队列并取回下一批任务，无需等待下游确认。跨硬件数据传输采用双缓冲机制——每个stage输出端维护写入缓冲与传输缓冲，写入完成后角色互换，DMA传输与下一请求的计算完全重叠。具体传输路径：CPU$\rightarrow$310P走PCIe 4.0（文本嵌入$\sim$64KB，延迟约2$\mu$s），310P$\rightarrow$910C走HCCS（编码特征$\sim$100MB，延迟约2ms），910C$\rightarrow$910B走HCCS（潜空间张量$\sim$8MB，延迟约0.16ms）。

系统采用流水线滑窗控制并发请求数上限$K_{\text{max}}=16$，按MMDiT批大小上限（4）与各stage处理延迟比（$\approx$1:2:50:10）归一化分配窗口配额，将峰值内存控制在理论最大值的65\%以内。
\subsection{分层启动缓存}

\subsubsection{三级缓存结构}

构建与stage分离架构深度适配的三级分层缓存：

\begin{enumerate}
    \item \textbf{模型权重缓存}：缓存各stage对应的模型主干FP16权重。文本预处理$\sim$168MB（84M参数），LLM编码$\sim$14GB（7B参数，加载至310P 24GB HBM），MMDiT$\sim$40GB（20B参数，加载至910C 64GB HBM），VAE$\sim$200MB（100M参数，加载至910B 32GB HBM）。权重按stage粒度独立加载与卸载，四设备并行加载将冷启动时间从全量串行加载的$\sim$60s压缩至$\sim$12s。
    \item \textbf{中间特征缓存}：缓存前序stage的输出特征张量，避免重复执行上游推理。核心条目包括：（a）文本编码特征——提示词经LLM编码后的条件向量，OpenPrompt等预置风格模板命中率可达90\%以上；（b）VAE图像编码特征——输入图像的潜空间表示，复用于相同图像的多次编辑；（c）风格LoRA嵌入——高频风格权重驻留在NPU HBM中。特征缓存采用主机内存（DDR）+设备内存（HBM）两级存储层次。

    在MMDiT Diffusion stage内部，本系统进一步引入CacheDiT~\cite{CacheDiTTrainingfreeCache2025}的块级特征缓存技术以加速逐step去噪过程。具体集成三项子策略：（a）\textbf{DBCache（Dual Block Cache）}：缓存在相邻去噪步间输出变化趋于稳定的前$n$个Transformer block（F$n$）和末$n$个block（B$n$）的中间激活值，在后续step中复用缓存值跳过这些block的重计算，配置为F8B8时可将去噪步计算量降低约40\%；（b）\textbf{FBCache（First Block Cache）}：轻量级变体，仅缓存第一个block的输出，适用于batch内负载不均时的选择性加速，计算开销约为DBCache的1/8；（c）\textbf{Cache CFG}：在分类器自由引导（Classifier-Free Guidance）推理中，将无条件分支的预测结果单独缓存，相邻step间无条件分支的特征漂移小于条件分支，实现额外1.3—2.0$\times$的CFG加速。
    \item \textbf{推理参数缓存}：缓存采样器类型、CFG scale、随机种子、去噪步数、输出分辨率等运行时参数，以JSON结构化按会话ID索引存储于CPU内存。FlashPS~\cite{FlashPSEfficientGenerativeJiang2026}在图像编辑场景中采用了按步状态缓存的思路，本工作将其扩展至通用文生图全流程的多级联合缓存。
\end{enumerate}

\subsubsection{缓存预加载与淘汰策略}

系统启动时，四个设备的权重并行加载（CPU$<$1s，310P$\sim$5s，910C$\sim$12s，910B$<$1s）。中间特征缓存在系统初始化后根据历史trace预计算Top-200高频提示词模板的文本编码特征（$\sim$3s），CacheDiT的FBCache在首个请求的前几步去噪中自动完成块级特征预热。

淘汰策略上，权重缓存采用基于stage滑动窗口利用率的LRU：当某设备剩余显存低于10\%时卸载利用率最低的非活跃stage权重至主机内存，重载速度约为初始加载的3—5倍。中间特征缓存沿用LMCache~\cite{LMCacheEfficientKVLiu2025}的层级迁移思路：设备内存中低频（访问计数$f<3$）且最久未用的条目迁移至主机内存，主机内存中超过TTL（$T_{\text{TTL}}=30$min）的条目彻底删除，维持HBM命中率$\ge92\%$。推理参数缓存采用会话级TTL（1h）惰性淘汰。
\subsection{结果聚合模块}
-
\section{实验评估}

\subsection{实验设置}

\paragraph{硬件环境} 
CPU：Kunpeng-920；昇腾NPU-310P（显存24GB）；昇腾NPU-910C（显存64GB）；昇腾NPU-910B（显存32GB）。

% ===== 5. 结论 =====
\section{结论}
-

% 参考文献左对齐
\printbibliography[
  heading=bibliography,
  title={\zihao{5}\simkai\textbf{参考文献}}
]

\end{document}
